\section{Verification method and acceptance criteria}
An inference result is evaluated against the reference appropriate to the claim. Exact comparison uses an independently implemented integer model. The floating-point result describes the effect of hardware representation and approximation. The transmitted target bits address decoding quality.

\begin{diagram}{7mm}
\node[port,text width=29mm] (i) {Prepared inputs\\and parameters};
\node[box,text width=33mm,right=13mm of i,yshift=8mm] (f) {Integer reference\\expected words};
\node[core,text width=33mm,right=13mm of i,yshift=-8mm] (r) {Lowered RTL\\observed words};
\node[port,text width=31mm,right=13mm of f,yshift=-8mm] (c) {Compare values\\and positions};
\draw[edge] (i.east)--(f.west);\draw[edge] (i.east)--(r.west);
\draw[edge] (f.east)--(c.west);\draw[edge] (r.east)--(c.west);
\end{diagram}
\capof{fig:verify}{Exact comparison for the trained frame. Both paths consume the same quantized image and prepared input. They calculate their outputs independently.}

{\smalltbl
\begin{tabularx}{\linewidth}{@{}P{40mm}YY@{}}
\toprule
Reference or test & Acceptance question & Example evidence\\ \midrule
Integer reference and RTL & Are the observed words equal to independently computed words? & Exact logit equality, index sequence, termination\\
Intermediate comparison & Did the right address, format and value commit? & Fixture checks of computational scratch writes\\
PyTorch and integer model & How different are the two representations on a sample? & Absolute score difference and decision agreement\\
Transmitted target data & Does decoding recover the original data? & Bit-error rate on a representative distribution\\
Protocol and arithmetic & Do local transactions and operations meet the contract? & Directed cases and scoped proofs\\
\bottomrule
\end{tabularx}}

The sample-zero worked example uses a two-state CXXRTL model exported by Yosys from the SystemVerilog with the Slang frontend. The C++ harness loads parameters and stimulus, then checks output words against the independent reference \cite{labbundle}. This is a simulation of a lowered RTL design. Gate-level simulation and physical measurements would address additional questions. The simulator-independent SystemVerilog plan in the repository documents a proposed driver and monitor structure. Its planned cases are distinct from the executed evidence reported here \cite{rtlrepo}.
\clearpage

\section{Trained-frame results}
The worked example uses sample zero, with shape \trained. The run loaded \(687{,}991\) parameter words, then accepted prepared features and a padding mask for 144 encoded positions. In \code{comparison.csv}, all 144 position indices are distinct and all expected and observed signed integers match. The session reports successful completion \cite{labbundle}.

{\smalltbl
\begin{tabularx}{\linewidth}{@{}P{31mm}P{35mm}P{35mm}Y@{}}
\toprule
Position & Expected F10 word & Observed RTL word & Result\\ \midrule
0 & \(-680\) & \(-680\) & Exact match\\
1 & \(2781\) & \(2781\) & Exact match\\
2 & \(1929\) & \(1929\) & Exact match\\ \midrule
All 144 positions & \multicolumn{2}{l}{144 distinct indices, zero mismatches} & Exact agreement\\
\bottomrule
\end{tabularx}}
\source{Trained shape \trained, sample zero, one read slot. See \code{comparison.csv}, \code{frame.json}, and \code{run.log} in the lab bundle \cite{labbundle}.}

\subsection{Example of the two comparisons}
At position zero the floating-point PyTorch logit is \(-0.66001004\). The fixed-point model returns the signed integer \(-680\), corresponding to
\[
 \frac{-680}{2^{10}}=-0.6640625.
\]
The difference between the PyTorch value and the interpreted fixed-point value reflects the two numerical models. The exact RTL result is the equality of expected and observed integers at \(-680\). Output-word agreement and floating-point agreement are therefore separate comparisons.

The 144 entries of the exported score comparison have mean absolute PyTorch-to-fixed-point logit difference \(0.01956\) and maximum \(0.06853\). The PyTorch and fixed-point decisions agree at all 144 positions in this sample. Decision agreement between representations differs from comparing decisions to transmitted target bits, and gives no estimate of a population bit-error rate.

\aside{Result scope}{The trained run uses one read slot and clear padding flags. It validates the observed output words for this frame. Multiread aggregation and padded-query behavior require other evidence, such as directed fixtures. The result provides no timing, power, or throughput measurement.}
\clearpage
\section{Ten-frame SystemVerilog simulation}
The lab ran the full trained configuration \trained{} using the SystemVerilog testbench \code{trained_frame_tb.sv}. It loaded the exported parameter image, provided the prepared features and mask, and checked the 144 returned indexed logits against the independent fixed-point Python outputs for each frame \cite{backend}.

{\smalltbl
\begin{tabularx}{\linewidth}{@{}P{32mm}YP{34mm}P{27mm}@{}}
\toprule
Frames & Simulator & Words compared & Mismatches\\ \midrule
00 & Cadence Xcelium & 144 & 0\\
01--09 & Synopsys VCS & 1,296 & 0\\ \midrule
All ten & & 1,440 & 0\\
\bottomrule
\end{tabularx}}

The logs report \(1{,}533{,}280{,}239\) cycles per frame and status \code{0001}, which records rounding. A separate comparison between fixed-point and floating-point decoded bits found nine differences near zero across the ten-frame set. This comparison measures the effect of numerical representation. The integer RTL check matched every output word. The measured cycle count belongs to the trained configuration \cite{backend}.
\clearpage

\section{Regression evidence and reproduction}
The corrected RTL regression recorded on 17 October 2026 adds coverage around arithmetic, normalization, protocol errors, and intermediate tensor writes. These tests use small configurations and controlled stimuli. Their purpose is to isolate issues that could remain hidden at the final output \cite{rtlrepo}.

{\smalltbl
\begin{tabularx}{\linewidth}{@{}P{40mm}P{33mm}Y@{}}
\toprule
Check & Recorded coverage & What it checks\\ \midrule
Complete-core execution & 16 fixture frames in four configurations & Stage progress and \(2{,}705\) intermediate writes\\
Arithmetic services & 958 vectors & Values, formats, flags, and response behavior\\
LayerNorm & 53 cases & Edge statistics, ties, and selected precision cases\\
Controller errors and resets & 22 directed scenarios & Delayed responses, reset, stalls, and recovery\\
Connected faults & Four cases & Integrated failure and recovery behavior\\
Conversions and memory & Exhaustive conversion and bounded properties & Defined rounding, saturation, format and routing cases\\
\bottomrule
\end{tabularx}}
\source{Corrected implementation record \code{RTL-REVIEW-FIXES-2026-10-17.md} \cite{rtlrepo}. This campaign is separate from the earlier 16 October baseline and the trained-frame CXXRTL run.}

\subsection{Reproducing the checks}
\begin{enumerate}
\setlength{\itemsep}{4pt}
\item Start at the source revision cited below. Set up the pinned repository requirements and confirm the top-level dimensions and interface revision. Run \code{python scripts/check.py smoke} for elaboration and manifest integration.
\item Use \code{python scripts/check.py numerics} for the numerical sweeps, \code{regressions} for directed and proof cases, and \code{full} for the complete-core group. The verification guide defines the scope of each and the output logs.
\item For the trained case, obtain the local checkpoint, effective configuration, features, mask, expected outputs, and parameters from the project handoff. Match their provenance first. Compare order and raw words, then inspect the simulation log and expected-output CSV.
\end{enumerate}

The repository checks and lab run have different environments. The source tree contains the architecture, RTL, and fixture drivers. The trained weights and raw lab export are separate project materials. The numbers above are recorded results from the cited runs.
\clearpage

\section{Logic synthesis and mapped cell area}
The implementation exercise used the \code{bcjrformer_dc_top} wrapper and shape \smallcore. The smaller case allows Design Compiler to exercise elaboration, library mapping and constraints independently of the trained-model run. The comparison uses Synopsys Design Compiler U-2022.12 and the Tower FS120 180\,nm typical standard-cell library named \code{tsl18fs120_typ}. The RTL weight arrays remained reloadable through the host interface. Design Compiler mapped the arrays into standard cells and reported zero SRAM macros \cite{labbundle,backend}.

{\smalltbl
\begin{tabularx}{\linewidth}{@{}P{34mm}P{42mm}Y@{}}
\toprule
Synthesis constraint & Reported mapped cell area & Available evidence\\ \midrule
20\,ns period & \(1{,}174{,}109.50\) report units & Direct \code{bcjr-small-area.rpt}\\
200\,ns period & \(1{,}052{,}431.25\) report units & Area value printed in \code{dc-session.log}\\
\bottomrule
\end{tabularx}}
\source{Shape \smallcore, Tower 180\,nm, typical standard-cell library, Design Compiler synthesis. The 200\,ns primary report is named in the session log. The value here is transcribed from that log \cite{labbundle}.}

\begin{diagram}{1mm}
\fill[Fog] (0,1.08) rectangle (11.5,1.50);
\fill[DeepTeal] (0,0.09) rectangle (10.31,0.51);
\node[anchor=west,font=\small] at (0,1.77) {20 ns target};
\node[anchor=west,font=\small] at (0,0.80) {200 ns target};
\end{diagram}
\capof{fig:area}{Relative mapped cell area across the two documented runs. Both bars share a common scale.}

The 200\,ns run reports about \(10.4\%\) less mapped cell area than the 20\,ns baseline. The comparison reflects two runs and does not isolate the effect of the clock target from other mapping choices. The 20\,ns report lists the normalization engine at \(30.2\%\) and the activation service at \(27.4\%\) of mapped cell area. These baseline shares identify two places to examine when investigating arithmetic cost.

The baseline report flags a capacitance-unit conflict and a high-fanout estimate. Mapped-cell area remains in report units. The reported net-interconnect area is a separate modeled quantity. The physical results below report the core boundary separately.
\clearpage

\section{Synthesis timing}
Synthesis static timing compares a path's arrival time with its required time. For setup, slack is required time minus arrival time. A negative value means the path misses the constraint. Separate hold analysis checks the minimum delay. The timing values depend on the library, constraint, design state, and analysis stage.

{\smalltbl
\begin{tabularx}{\linewidth}{@{}P{34mm}P{28mm}P{27mm}Y@{}}
\toprule
Clock constraint & Setup slack & Hold slack & Available evidence\\ \midrule
20\,ns / 50\,MHz target & \(-145.64\)\,ns & Unavailable & Direct timing report and SDC\\
200\,ns / 5\,MHz target & \(+0.05\)\,ns & \(+0.32\)\,ns & Setup and hold excerpts printed in the session log\\
\bottomrule
\end{tabularx}}
\source{Small-core shape \smallcore, Tower 180\,nm typical standard-cell library, Design Compiler synthesis. The 20\,ns report uses a wire-load model. The 200\,ns values come from \code{dc-session.log} \cite{labbundle}.}

The worst reported setup path at 20\,ns lies in the LayerNorm logic. Its arrival time is \(165.52\)\,ns and its required time is \(19.88\)\,ns, giving the reported \(-145.64\)\,ns slack. The separate 200\,ns run reports positive setup and hold slack, with \(50\)\,ps of setup margin. The frequency equivalent of the 200\,ns constraint is 5\,MHz. The post-route analysis below uses the small core. The trained-frame cycle count belongs to a different configuration.

\clearpage

\section{Small-core physical implementation}
Cadence Innovus 23.33 implemented \code{bcjrformer_dc_top} on the Tower 180\,nm six-metal stack. The mapped netlist and adapted constraints were imported at a 200\,ns trial period. The core was approximately \(5.142 \times 5.135\)\,mm with 917 placement rows and about 50\% initial utilization. Initial placement reported 730,967 cells placed and zero unplaced \cite{backend}.

\begin{enumerate}
\item \textbf{Power and pins.} VDD/VSS rings used horizontal M3 and vertical M4, with special routing connecting the standard-cell rails. The 142 fixed core boundary pins grouped the host weight interface on the west, frame inputs on the south, results on the east, and clock and reset on the north.
\item \textbf{Placement and clocks.} Pre-CTS optimization improved worst setup slack from \(-19.942\) to \(-0.402\)\,ns in the trial setup. Clock-tree synthesis built a tree for 8,801 sinks. The early timing results preceded corrections to clock propagation and post-route clock RC scaling.
\item \textbf{Routing and fillers.} Global and detail routing initially produced 94 antenna violations. Router diode insertion reduced the count to three, and three manual diode attachments cleared the remainder. The lab filler flow inserted 1,304,530 cells. The layout reported 46.675\% functional placement density and 100\% occupancy including fillers.
\end{enumerate}

\section{Post-route timing and slew correction}
The final analysis used propagated clocks, unity post-route clock RC scale factors, and on-chip variation with common path pessimism removal. The lab extraction setup used the six-metal capacitance table. The recorded extraction settings and warnings still need review for foundry sign-off \cite{backend}.

At 200\,ns, the corrected slow-corner analysis reported \(-67.570\)\,ns setup WNS, approximately \(-29{,}915.4\)\,ns TNS, and 478 violating paths. The worst path started at the LayerNorm variance register and went through the reciprocal square root. The lab adopted a 300\,ns target, approximately 3.33\,MHz.

\clearpage
\subsection{Final checks at 300\,ns}
The initial 300\,ns setup run passed with four rising-edge transition violations. An engineering change upsized one NOR and three inverters. The larger cells initially caused eight placement overlaps in the filler-filled core. Removing 340 local fillers cleared the overlaps. Incremental routing, refill, and verification followed \cite{backend}.

{\smalltbl
\begin{tabularx}{\linewidth}{@{}P{24mm}YP{22mm}P{22mm}P{20mm}@{}}
\toprule
Check & Corner & WNS (ns) & TNS (ns) & Failing paths\\ \midrule
Setup & Slow, 1.65\,V, \(125^{\circ}\mathrm{C}\) & \(+27.961\) & 0.000 & 0\\
Hold & Fast, 2.017\,V, \(-40^{\circ}\mathrm{C}\) & \(+0.141\) & 0.000 & 0\\
\bottomrule
\end{tabularx}}
\source{Both final reports cover 8,836 paths in the small core. See the slow setup and fast hold reports documented in the backend handoff \cite{backend}.}

The slow setup report showed zero max-transition, max-capacitance, max-fanout, and max-length violations. Innovus geometry and antenna checks found zero violations. Signal and power connectivity checks found no problems or warnings.

\subsection{Exports and verification scope}
The refilled post-ECO layout was exported as \code{bcjrformer_dc_top.gds}, alongside the matching physical Verilog \code{bcjrformer_dc_top_lvs.v}. The layout is a routed core block with boundary pins. Foundry-deck DRC, extracted LVS, applicable metal-density checks, and gate-level equivalence remain pending. The book also has no verified power measurement. Physical implementation of the full trained configuration remains future work.
\clearpage
